% ============================================================
% CHAPTER 10 — Project 7: Full-Stack Integration
% ============================================================
\chapter{Project 7: Full-Stack Integration}\index{full-stack}\index{monorepo}

\section*{What You Will Build}
A fully integrated full-stack application:
\begin{itemize}
  \item Backend and frontend that communicate without errors
  \item CORS management and proxy configured correctly
  \item Unified end-to-end error handling
  \item One single project \texttt{\_CONTEXT.md} that governs both
  \item The Multi-Agent pattern applied in practice
  \item A reproducible development system
\end{itemize}

\textbf{Estimated time}: 45--60 minutes\\
\textbf{Prerequisite}: Working Backend (Ch.~6--8) and Frontend (Ch.~9)

% ---------------------------------------------------------
\section{The Monorepo Structure}

\begin{praticabox}[title={Hands-On --- Organize the workspace}]
\begin{lstlisting}[language={}]
notes-fullstack/
+-- _CONTEXT.md              <- Global context
+-- SKILL.md                 <- Shared skill
+-- PROGRESS.md              <- Persistent memory
+-- backend/
|   +-- _CONTEXT.md          <- Backend context
|   +-- package.json
|   +-- prisma/
|   +-- src/
+-- frontend/
    +-- _CONTEXT.md          <- Frontend context
    +-- SKILL.md             <- React skill
    +-- package.json
    +-- src/
\end{lstlisting}
\begin{lstlisting}[language=bash]
mkdir notes-fullstack
mv notes-api notes-fullstack/backend
mv notes-frontend notes-fullstack/frontend
cd notes-fullstack && git init
\end{lstlisting}
\end{praticabox}

% ---------------------------------------------------------
\section{The Global Context}

The \texttt{\_CONTEXT.md} in the root is the most important document --- it describes the entire system.

\begin{lstlisting}[language={},title={root \texttt{\_CONTEXT.md}}]
# Project: Notes Full-Stack Application

## Overview
Full-stack web application for managing personal notes
with OAuth 2.0 authentication.

## Architecture
\end{lstlisting}

\begin{figure}[H]
\centering
\includegraphics[width=0.85\textwidth]{10-fullstack-dev-arch}
\caption{Full-Stack Architecture in Development}
\end{figure}

\begin{lstlisting}[language={}]
## API Contract - Format
Success: { "success": true, "data": T }
Error:   { "success": false,
            "error": { "message": "...", "code": "..." } }

## Error codes
- 401: Not authenticated -> redirect to login
- 403: Not authorized -> show error
- 404: Resource not found -> 404 page
- 422: Validation failed -> field errors
- 500: Server error -> generic message
\end{lstlisting}

\begin{notebox}
This document is the contract between the two parts of the application. When you modify the backend, the AI knows what to expect from the frontend and vice versa.
\end{notebox}

% ---------------------------------------------------------
\section{CORS Management in Development and Production}\index{CORS}

\textbf{In development}: the Vite proxy (Chapter~9) forwards the requests.

\textbf{In production}: the backend directly serves the frontend static files.

\begin{praticabox}[title={Hands-On --- Backend for production}]
\begin{lstlisting}[language={}]
Modify the backend to serve the frontend static files
in production.

In app.js:
- If NODE_ENV === 'production', serve files from
  ../frontend/dist
- For non-API routes, serve index.html (React Router)
- CORS only in development
\end{lstlisting}
\end{praticabox}

% ---------------------------------------------------------
\section{End-to-End Error Handling}

\begin{figure}[H]
\centering
\includegraphics[width=0.85\textwidth]{09-error-flow}
\caption{Error Flow}
\end{figure}

\begin{praticabox}[title={Hands-On --- Complete error handling}]
\begin{lstlisting}[language={}]
Implement end-to-end error handling:

1. Backend: global middleware that converts all errors
   to the standard format. Handle PrismaClientKnownRequest
   (P2025 = not found), ValidationError, AuthError.

2. Frontend: Axios interceptor that:
   - 401: clears auth, redirects to login
   - 422: returns validation errors to the form
   - 500: generic error toast
   - Network error: banner "Connection lost"

3. Frontend: React ErrorBoundary for rendering errors.
\end{lstlisting}
\end{praticabox}

\begin{warnbox}
The backend error handler \textsc{must not} ever expose internal details in production. Messages such as ``Cannot read property `id' of undefined'' or stack traces are security vulnerabilities (Information Disclosure, OWASP~A01~\cite{owasp:top10}).
\end{warnbox}

% ---------------------------------------------------------
\section{Development Session with \texttt{PROGRESS.md}}

\begin{praticabox}[title={Hands-On --- Create \texttt{PROGRESS.md}}]
\begin{lstlisting}[language={},title={Example progress file}]
# Notes Full-Stack -- Progress

## Current Status
- [x] Backend REST API with Express.js
- [x] PostgreSQL database with Prisma
- [x] OAuth 2.0 authentication
- [x] React frontend with notes dashboard
- [x] Complete notes CRUD (frontend <-> backend)
- [x] End-to-end error handling
- [ ] Automated testing (Chapter 14)
- [ ] Production deployment (Chapter 15)

## Architectural Decisions
- JWT in httpOnly cookie (not localStorage)
- Vite proxy in dev, static serving in production
- Monorepo with separate contexts
- Context API for auth state (no Redux)

## Resolved Problems
- CORS: Vite proxy + conditional cors middleware
- Refresh token: Axios interceptor retries the request
- OAuth redirect: callback -> backend -> frontend
\end{lstlisting}
\end{praticabox}

% ---------------------------------------------------------
\section{The Multi-Agent Pattern in Practice}

It is not about making agents communicate with each other, but about \textbf{changing the role of the AI} in the different phases of the work.

\subsection{The Planner-Generator-Evaluator Pattern}

\textbf{1.\ Planner:}
\begin{lstlisting}[language={}]
Analyze the root _CONTEXT.md. I want to add
"categories" to notes. DO NOT write code. List:
1. Which backend files must be modified and why
2. Which frontend files must be modified and why
3. The database migrations required
4. The recommended implementation order
\end{lstlisting}

\textbf{2.\ Generator:}
\begin{lstlisting}[language={}]
Implement point 1 of the plan: the database migration.
Add the Category model and the relation with Note.
\end{lstlisting}

\textbf{3.\ Evaluator:}
\begin{lstlisting}[language={}]
Review the code you just generated. Verify:
1. Is the migration reversible?
2. Do the endpoints follow the standard format?
3. Does the frontend handle the "no category" case?
4. Can one user access another user's categories?
List any problems found.
\end{lstlisting}

\begin{tipbox}
Asking ``show me a diff'' before applying the changes is an effective defensive practice.
\end{tipbox}

% ---------------------------------------------------------
\section{Development Scripts}

\begin{praticabox}[title={Hands-On --- Root npm scripts}]
\begin{lstlisting}[language=json,title={root \texttt{package.json}}]
{
  "name": "notes-fullstack",
  "private": true,
  "scripts": {
    "dev": "concurrently \"npm run dev:backend\"
            \"npm run dev:frontend\"",
    "dev:backend": "cd backend && npm run dev",
    "dev:frontend": "cd frontend && npm run dev",
    "build": "cd frontend && npm run build",
    "db:migrate": "cd backend && npx prisma migrate dev",
    "db:studio": "cd backend && npx prisma studio"
  },
  "devDependencies": {
    "concurrently": "^8.0.0"
  }
}
\end{lstlisting}
\begin{lstlisting}[language=bash]
npm install
npm run dev  # Start backend and frontend together
\end{lstlisting}
\end{praticabox}

% ---------------------------------------------------------
\section{End-to-End Testing}

\begin{center}
\begin{tabularx}{\textwidth}{l X}
\toprule
\textbf{Test} & \textbf{Verification} \\
\midrule
Cold start & \texttt{npm run dev} $\rightarrow$ everything starts without errors \\
Google login & Click $\rightarrow$ OAuth $\rightarrow$ Dashboard \\
Complete CRUD & Create $\rightarrow$ Read $\rightarrow$ Edit $\rightarrow$ Delete \\
Logout & Logout $\rightarrow$ Try /dashboard $\rightarrow$ Redirect \\
Expired token & The app renews automatically \\
Backend down & Frontend shows ``Connection lost'' \\
Mobile & Responsive layout on emulator \\
\bottomrule
\end{tabularx}
\end{center}

\begin{checkpointbox}
If all tests pass, the full-stack application is complete.
\end{checkpointbox}

% ---------------------------------------------------------
\section*{What You Have Built So Far}

Take a moment. Starting from zero, in 10 chapters you have built:

\begin{enumerate}
  \item An \textbf{interactive} Python program (Ch.~4)
  \item A \textbf{task manager CLI} with persistence (Ch.~5)
  \item A \textbf{REST API} with Express.js and Swagger (Ch.~6)
  \item A \textbf{PostgreSQL database} with Prisma ORM (Ch.~7)
  \item An \textbf{OAuth~2.0 authentication} system (Ch.~8)
  \item A \textbf{React frontend} with dashboard and CRUD (Ch.~9)
  \item A \textbf{fully integrated full-stack application} (Ch.~10)
\end{enumerate}

All \textbf{without writing a single line of code manually}.

% ---------------------------------------------------------
\section*{Summary}

\begin{center}
\begin{tabularx}{\textwidth}{l X}
\toprule
\textbf{Aspect} & \textbf{Detail} \\
\midrule
Architecture & Monorepo with separate contexts \\
Global context & root \texttt{\_CONTEXT.md} with API contract \\
CORS & Vite proxy (dev) + static serving (prod) \\
Error handling & End-to-end from the database to the UI toast \\
Multi-Agent & Planner-Generator-Evaluator pattern \\
Memory & \texttt{PROGRESS.md} for ongoing sessions \\
Scripts & \texttt{npm run dev} starts everything with concurrently \\
\bottomrule
\end{tabularx}
\end{center}
